% Generated draft: requires mathematical and source review.
\documentclass[11pt]{article}
\usepackage[margin=25mm]{geometry}
\usepackage{fontspec}
\setmainfont{Latin Modern Roman}
\usepackage{amsmath,amssymb,mathtools,unicode-math}
\setmathfont{Latin Modern Math}
\usepackage{enumitem,needspace,xurl,hyperref}
\setlength{\parindent}{0pt}
\setlength{\parskip}{5pt}
\setlength{\emergencystretch}{3em}
\newcommand{\N}{\mathbb{N}}
\newcommand{\Z}{\mathbb{Z}}
\newcommand{\Q}{\mathbb{Q}}
\newcommand{\R}{\mathbb{R}}
\newcommand{\C}{\mathbb{C}}
\newcommand{\F}{\mathbb{F}}
\newcommand{\A}{\mathbb{A}}
\newcommand{\PP}{\mathbb{P}}
\newcommand{\E}{\mathbb{E}}
\newcommand{\eps}{\varepsilon}
\newcommand{\dd}{\,\mathrm{d}}
\begin{document}
\section*{UnsolvedMath ID: 30003077}
\textbf{Source problem code:} catalogue entry
\subsection{Definitions and mathematical statement}
The supplied problem does not specify the ambient space, the ground field, or whether the linear system is complete. Consequently, its numerical target can be expressed precisely only relative to these additional data.

For any specified geometric setting in which curves of degree \(d\) are represented by sections of a line bundle \(L_d\) on an ambient space \(X\) over a field \(k\), let \(Z=\{Q_1,\ldots,Q_s\}\) be the prescribed finite set. If multiplicities are prescribed, write \(a_j\geq 1\) for the multiplicity at \(Q_j\); otherwise put \(a_j=1\). Define the base subscheme \(B\) by
\[
\mathcal I_B=\bigcap_{j=1}^{s}\mathcal I_{Q_j}^{a_j}.
\]
Represent the given linear system by \(\mathbb P(V)\), where \(V\) is a specified finite-dimensional \(k\)-linear subspace of \(H^0(X,L_d\otimes\mathcal I_B)\). The complete system corresponds to equality, but completeness is not stated in the source.

For a specified number \(r\geq 1\) and specified positive integers \(m_1,\ldots,m_r\), a collection of fat points supported at distinct points \(P_1,\ldots,P_r\) is the subscheme \(F\) with
\[
\mathcal I_F=\bigcap_{i=1}^{r}\mathcal I_{P_i}^{m_i}.
\]
The number of independent conditions that \(F\) imposes on \(V\) is
\[
 c_V(F)=\operatorname{rank}_k\!\left(V\longrightarrow H^0(F,L_d|_F)\right)
 =\dim_k V-\dim_k\!\left(V\cap H^0(X,L_d\otimes\mathcal I_F)\right).
\]
In particular, imposing both sets of conditions uses \(\mathcal I_B\cap\mathcal I_F\), including when their supports overlap.

\textbf{Problem, conditional on specifying the missing geometric data.} Determine \(c_V(F)\) for general choices of the supports \(P_i\), as a function of the degree, the prescribed base conditions, the chosen linear system, and the fat-point multiplicities. Here a general choice means a choice in a nonempty Zariski-open subset of the specified irreducible parameter space of allowed configurations. The intended parameter space, any restriction excluding \(Z\), and the permitted geometric data must be supplied before this becomes a unique classification problem. This formulation records the requested rank; it does not assert that the rank equals an expected dimension.
\subsection{Short English statement}
Determine how many independent linear conditions general fat points impose on a fixed-degree system of curves already passing through a prescribed finite set \(Z\), possibly with specified multiplicities at \(Z\). The ambient geometry, field, admissible configurations, and precise linear system remain unspecified in the supplied record.
\subsection{Sources}
\begin{itemize}
\item Local UnsolvedMath registry, supplied record with permanent ID \(30003077\), title \emph{Conditions Imposed by Fat Points on Curve Systems}, exact source problem code \texttt{catalogue entry}. The record assigns the problem to the Oberwolfach Reports Open Problems collection and gives proposed year \(2016\); no individual proposer is supplied.
\item Supplied source citation: \emph{Mini-Workshop: Arrangements of Subvarieties, and their Applications in Algebraic Geometry} (2016), Oberwolfach Reports, pp.\ 653--702, DOI \texttt{10.4171/owr/2016/14}. Source URL and literature-source link labelled ``OWR report 14/2016'': \url{https://doi.org/10.4171/owr/2016/14}. The supplied background describes this report as documentation of a mathematical workshop and its open-problem discussions.
\item Supplied literature-source attribution: Cook--Harbourne--Migliore--Nagel, \url{https://doi.org/10.1112/S0010437X18007376}. No title, specific result, or pinpoint reference is supplied for this item.
\item Literature assessment dated exactly \texttt{2026-08-21}: the supplied curation labels the problem open and describes it as a generalized SHGH-type classification problem for linear subsystems with imposed base points. It states that the classical SHGH problem remains unresolved and that the broader proposal is not settled. This is supplied dated curation derived from the OWR record, not an independently verified literature assessment; time-sensitive status and cited claims require independent verification.
\item The registry labels the statement \texttt{corrected\_verified}, with the supplied explanation that notation, typography, and sentence structure were repaired without changing mathematical content. This records an editorial qualification and does not establish mathematical correctness.
\end{itemize}
\end{document}
